\documentclass[11pt]{article}
\usepackage{mathblatt}
% Präambel des Setzers (werkzeuge/setzer-praeambel.tex, wortgleich); dazu \gteil (Muster T2U)
\makeatletter
\setlength{\mbpfbe}{0mm}% keine Punkte (pruefung.md § 5: nur beim Prüfstein)
\newlength{\gEin}\setlength{\gEin}{0pt}
\let\gSchrift\relax
\renewcommand{\pfkreuzzeile}[1]{\par\noindent\hangindent3.2em\hangafter1\hspace*{1.6em}$\square$\ #1\par}
\newcommand{\gkreuz}[1]{$\square$\ #1\hspace{2em}}
\newcommand{\gkreuzzeile}[2]{\par\noindent\hangindent2.2em\hangafter1\mbox{#1\,$\square$}\ #2\par}% Buchstabe, Option: eine je Zeile (neben Skizze)
% links, was man ansieht, rechts, was man tut; Spaltengrenze überall an derselben Stelle (gemeinsam.md „Satz“)
\newsavebox{\gbox}\newlength{\gLinks}
\newcommand{\gzwei}[2]{\par\smallskip\noindent\sbox{\gbox}{#1}%
  \setlength{\gLinks}{\dimexpr0.5\textwidth-\gEin-\mbpfnr\relax}%
  \ifdim\wd\gbox>\dimexpr\gLinks-4mm\relax
    \usebox{\gbox}\par\smallskip\noindent\begin{minipage}{\linewidth}#2\end{minipage}%
  \else
    \begin{minipage}[t]{\gLinks}\vspace{0pt}\usebox{\gbox}\end{minipage}%
    \begin{minipage}[t]{\dimexpr\linewidth-\gLinks\relax}\vspace{0pt}\raggedright #2\end{minipage}%
  \fi\par}
% v0.9 (12.10., pruefung.md § 5 „Marke mit B, Bedingungen rechts“): drittes Argument von pfaufg = Bedingungen
% der Prüfung („Basis · ohne TR“, „FOR“), klein grau im rechten Rand auf der Grundlinie der Nummer, Zeilen mit \\
\newcommand{\gbedrand}[2]{\ifblank{#2}{}{\rlap{\hspace*{\dimexpr#1+2mm\relax}%
  \parbox[t]{15mm}{\raggedright\scriptsize\color{mbgrau}\strut#2}}}}
% Teil eines Bündels (a, b …): Marke links im Rand, Bedingungen rechts im Rand, auf der Zeile des Teils
\newcommand{\gteilrand}[2]{\llap{\mbpfrandmarke{#1}\hspace*{\mbpfnr}}\gbedrand{\linewidth}{#2}}
% Teil eines Bündels: Buchstabe in eigener Spalte, Inhalt eingerückt (\gEin = Einzug, damit die Spaltengrenze von
% \gzwei an derselben Stelle bleibt)
\newenvironment{gbteil}[3]{\par\smallskip\noindent\gteilrand{#2}{#3}\makebox[\mbpfteil][l]{\textbf{#1}}%
  \begin{minipage}[t]{\dimexpr\linewidth-\mbpfteil\relax}\setlength{\gEin}{\mbpfteil}\raggedright\strut\ignorespaces}%
 {\end{minipage}\par}
% Grafik höchstens so breit wie die Spalte (Bündel nebeneinander, Koordinatensystem ohne Karo)
\newsavebox{\gpbox}
\newcommand{\gpassend}[1]{\sbox{\gpbox}{#1}\ifdim\wd\gpbox>\linewidth\resizebox{\linewidth}{!}{\usebox{\gpbox}}\else\usebox{\gpbox}\fi}
\RenewDocumentEnvironment{pfaufg}{m m m}{%
  \begin{lrbox}{\mbaufgabebox}\begin{minipage}[t]{\linewidth}%
  \gSchrift\noindent\hspace*{\gEin}\mbpfrandmarke{#1}\gbedrand{\linewidth-\gEin}{#3}%
  \begin{minipage}[t]{\mbpfnr}\strut\textbf{#2}\end{minipage}%
  \begin{minipage}[t]{\dimexpr\linewidth-\gEin-\mbpfnr\relax}\raggedright\strut\ignorespaces}%
 {\end{minipage}%
  \end{minipage}\end{lrbox}%
  \setlength{\mbaufgabehoehe}{\dimexpr\ht\mbaufgabebox+\dp\mbaufgabebox\relax}%
  \ifdefempty{\gkopfnach}{\par\addvspace{5pt}\Needspace*{\mbaufgabehoehe}}{%
    % Stufenkopf, der mit dieser (ersten) Aufgabe zusammenbleibt: Platz für beide verlangen, dann den Kopf setzen
    % (12.10., Quadratische: ein Kopf stand allein unten, die Aufgabe mit Koordinatensystem auf der nächsten Seite)
    \par\Needspace*{\dimexpr\mbaufgabehoehe+40pt\relax}\gkopfnach\gdef\gkopfnach{}\par\addvspace{5pt}}%
  \noindent\usebox{\mbaufgabebox}\par\addvspace{5pt}}
% Stufenkopf vormerken: \gkopfvor{<Kopf>} – die nächste pfaufg setzt ihn, sobald beide zusammen Platz haben
\def\gkopfnach{}
\newcommand{\gkopfvor}[1]{\gdef\gkopfnach{#1}}
% Teilaufgabe mit Skizze links in gzwei: Buchstabe in eigener Spalte, darunter Text und Skizze
\newcommand{\gteil}[2]{\makebox[\mbpfteil][l]{\textbf{#1}}\parbox[t]{\dimexpr0.5\textwidth-\gEin-\mbpfnr-5mm-\mbpfteil\relax}{\raggedright #2}}
% Stufenkopf der Vorlage (mathblatt.sty 2026-10-12: Linie, Abstand, 13 pt, \pfjahre) ohne ihren Vorhalt von
% 14 Zeilen: der Setzer fordert vor dem Kopf 0,16 Seitenhöhe (Kopf + erste Aufgabe), mehr nicht
\setlength{\mbpfstufevorhalt}{0pt}
\makeatother
\newcommand{\karo}[1]{\par\smallskip\noindent\begin{tikzpicture}
  \clip (0,0) rectangle ({\linewidth-1mm},{#1*10mm});
  \draw[black!16,line width=0.3pt,step=5mm] (0,0) grid ({\linewidth},{#1*10mm});
  \end{tikzpicture}\par}
\newcommand{\kk}[1]{$\square$\,#1\hspace{0.9em}}
\newcommand{\linie}[1]{\,{\color{black!45}\rule[-1pt]{#1}{0.4pt}}\,}
\newcommand{\pkt}{\textperiodcentered{}}
% v0.7 (12.10.): Schülerfrage unter dem Stufenkopf; Ankreuzoption als Bild (Buchstabe und Kästchen darunter)
\newcommand{\gfrage}[1]{\par\nobreak\noindent\hspace*{\mbpfnr}{\small\itshape #1}\par\nobreak\addvspace{2pt}}
\newlength{\gbildbreite}\setlength{\gbildbreite}{0.32\linewidth}% Setzer setzt sie je Reihe (vier Bilder: 0,24)
\newcommand{\gbildoption}[3]{\begin{minipage}[t][#3][b]{\gbildbreite}\centering #2\par\smallskip\mbox{#1\,$\square$}\end{minipage}}
% Lineare (12.10.): Bild mit Buchstabe ohne Kästchen (Diagramm zuordnen); Aussage mit „wahr / falsch“ zum Ankreuzen
\newcommand{\gbildname}[3]{\begin{minipage}[t][#3][b]{\gbildbreite}\centering #2\par\smallskip\mbox{\textbf{#1}}\end{minipage}}
\newcommand{\gwahrfalsch}[2]{\par\noindent\hangindent2.2em\hangafter1\makebox[2.2em][l]{(#1)}#2\hfill\mbox{wahr $\square$\quad falsch $\square$}\par}
\begin{document}

\pfheftstil
\pflaufkopf[P10 \pkt{} Prüfung]{Lineare Funktionen}{TAL2}\pfseitevon
\pfheftkopf[B = Basisteil, ohne Taschenrechner]{Lineare Funktionen}{}
\pfstufenzeile{Gerade zeichnen \pkt{} Gleichung ablesen oder aus Punkten aufstellen \pkt{} Nullstelle \pkt{} Punktprobe \pkt{} Gleichung aus Text}

% ===================== Stufe: Gerade zeichnen =====================
\Needspace*{0.16\textheight}\pfstufekopf{Gerade zeichnen}{immer \pkt{} \pfjahre{26 \pkt{} 25 \pkt{} 24 \pkt{} 23}{}}
% Schülerfrage der Stufe (Gliederung; Einstieg nur von dort, nie aus der Bank – 12.10.)
\begin{pfaufg}{}{1.}{}
Welche Funktionen sind linear? Kreuze an. Wähl eine lineare aus und gib ihr m und n an: $m = \underline{\hspace{10mm}}$, $n = \underline{\hspace{10mm}}$.
\par\smallskip
\par\noindent\makebox[\mbpfteil][l]{\textbf{a)}}$f(x) = 3 - 2x$\hfill\kk{linear}
\par\noindent\makebox[\mbpfteil][l]{\textbf{b)}}$f(x) = x^2 + 1$\hfill\kk{linear}
\par\noindent\makebox[\mbpfteil][l]{\textbf{c)}}$f(x) = 0{,}5x$\hfill\kk{linear}
\par\noindent\makebox[\mbpfteil][l]{\textbf{d)}}$f(x) = \tfrac{4}{x}$\hfill\kk{linear}
\par\fusshilfe{\mbox{1:~a) und c)}}
\end{pfaufg}
% 2025-OS-B1f-m1 (Original 2025-OS-B1f, Schwere I, rang 1; herausgelöste Fassung: keine)
\begin{pfaufg}{’25~\pkt{}~B1f}{2.}{Basis~\pkt{}\\ohne TR}
Kreuze das Bild an, das den Graphen einer linearen Funktion zeigt.
\par\smallskip\noindent\setlength{\gbildbreite}{\dimexpr0.96\linewidth/4\relax}\sbox{\gbox}{\vbox{\begin{tikzpicture}[x=1.1cm,y=1.1cm,line width=0.6pt]\draw[-{Stealth[length=1.5mm]}] (-0.8,0)--(0.95,0);\draw[-{Stealth[length=1.5mm]}] (0,-0.6)--(0,1.05);\draw[line width=0.8pt] (-0.8,0.94)--(0.9,-0.42);\end{tikzpicture}\par\begin{tikzpicture}[x=1.1cm,y=1.1cm,line width=0.6pt]\draw[-{Stealth[length=1.5mm]}] (-0.8,0)--(0.95,0);\draw[-{Stealth[length=1.5mm]}] (0,-0.6)--(0,1.05);\draw[line width=0.8pt,domain=0.1:0.9,samples=40,smooth] plot (\x,{0.09/\x});\end{tikzpicture}\par\begin{tikzpicture}[x=1.1cm,y=1.1cm,line width=0.6pt]\draw[-{Stealth[length=1.5mm]}] (-0.8,0)--(0.95,0);\draw[-{Stealth[length=1.5mm]}] (0,-0.6)--(0,1.05);\draw[line width=0.8pt,domain=-0.6:0.6,samples=40,smooth] plot (\x,{3*\x*\x-0.4});\end{tikzpicture}\par\begin{tikzpicture}[x=1.1cm,y=1.1cm,line width=0.6pt]\draw[-{Stealth[length=1.5mm]}] (-0.8,0)--(0.95,0);\draw[-{Stealth[length=1.5mm]}] (0,-0.6)--(0,1.05);\draw[line width=0.8pt] (0.1,-0.3)--(0.5,0.5)--(0.1,0.7);\end{tikzpicture}}}\gbildoption{A}{\begin{tikzpicture}[x=1.1cm,y=1.1cm,line width=0.6pt]\draw[-{Stealth[length=1.5mm]}] (-0.8,0)--(0.95,0);\draw[-{Stealth[length=1.5mm]}] (0,-0.6)--(0,1.05);\draw[line width=0.8pt] (-0.8,0.94)--(0.9,-0.42);\end{tikzpicture}}{\dimexpr\ht\gbox/4+\baselineskip\relax}\hfill\gbildoption{B}{\begin{tikzpicture}[x=1.1cm,y=1.1cm,line width=0.6pt]\draw[-{Stealth[length=1.5mm]}] (-0.8,0)--(0.95,0);\draw[-{Stealth[length=1.5mm]}] (0,-0.6)--(0,1.05);\draw[line width=0.8pt,domain=0.1:0.9,samples=40,smooth] plot (\x,{0.09/\x});\end{tikzpicture}}{\dimexpr\ht\gbox/4+\baselineskip\relax}\hfill\gbildoption{C}{\begin{tikzpicture}[x=1.1cm,y=1.1cm,line width=0.6pt]\draw[-{Stealth[length=1.5mm]}] (-0.8,0)--(0.95,0);\draw[-{Stealth[length=1.5mm]}] (0,-0.6)--(0,1.05);\draw[line width=0.8pt,domain=-0.6:0.6,samples=40,smooth] plot (\x,{3*\x*\x-0.4});\end{tikzpicture}}{\dimexpr\ht\gbox/4+\baselineskip\relax}\hfill\gbildoption{D}{\begin{tikzpicture}[x=1.1cm,y=1.1cm,line width=0.6pt]\draw[-{Stealth[length=1.5mm]}] (-0.8,0)--(0.95,0);\draw[-{Stealth[length=1.5mm]}] (0,-0.6)--(0,1.05);\draw[line width=0.8pt] (0.1,-0.3)--(0.5,0.5)--(0.1,0.7);\end{tikzpicture}}{\dimexpr\ht\gbox/4+\baselineskip\relax}\par
\fusshilfe{\mbox{2:~A}}
\end{pfaufg}
% 2024-OS-B1i-m1 (Original 2024-OS-B1i, Schwere I, rang 2; herausgelöste Fassung: keine)
\begin{pfaufg}{’24~\pkt{}~B1i}{3.}{Basis~\pkt{}\\ohne TR}
Gegeben ist die Gerade $y = 1{,}5x + 1$. Zeichne den Punkt ein, an dem sie die y-Achse trifft.
\par\smallskip\noindent\begin{ksys}[xmin=-2,xmax=2,ymin=-2,ymax=2,karo=0.5]\end{ksys}\par
\fusshilfe{\mbox{3:~Punkt (0|1)}}
\end{pfaufg}
% Bündel zeichnen: 2024-OS-K3a-m1, 2026-FOR-K5a-m1
\begin{pfaufg}{}{4.}{}
Zeichne den Graphen von $f$ in das Koordinatensystem.
% 2024-OS-K3a-m1 (Original 2024-OS-K3a, Schwere I, rang 3; herausgelöste Fassung: keine)
% Zweitthema ausgeblendet (Spalten <feld>_einheit, pruefheft-p.md v0.9)
% 2026-FOR-K5a-m1 (Original 2026-FOR-K5a, Schwere I, rang 4; herausgelöste Fassung: keine)
% Zweitthema ausgeblendet (Spalten <feld>_einheit, pruefheft-p.md v0.9)
\par\smallskip\noindent\begin{minipage}[t]{\dimexpr\linewidth/2-2mm\relax}\raggedright\textbf{a)}\hspace{0.4em}{\scriptsize\color{mbgrau}nach~P10 ’24~\pkt{}~3a}\par\smallskip
Die lineare Funktion $f$ hat die Gleichung $f(x) = 2x + 3$.
\par\smallskip\noindent\gpassend{\begin{ksys}[xmin=-2,xmax=2,ymin=-1,ymax=7,karo=0.35]\end{ksys}}\par
\end{minipage}\hfill
\begin{minipage}[t]{\dimexpr\linewidth/2-2mm\relax}\raggedright\textbf{b)}\hspace{0.4em}{\scriptsize\color{mbgrau}nach~P10 ’26~\pkt{}~5a \pkt{} FOR}\par\smallskip
Die lineare Funktion $f$ hat die Gleichung $f(x) = -3x + 1$.
\par\smallskip\noindent\gpassend{\begin{ksys}[xmin=-2,xmax=2,ymin=-3,ymax=5,karo=0.35]\end{ksys}}\par
\end{minipage}\par
\fusshilfe{\mbox{4:~a)~Gerade durch (0|3) und (1|5), steil steigend; b)~Gerade durch (0|1) und (1|$-$2), fallend}}
\end{pfaufg}
% 2023-OS-K4a-m1 (Original 2023-OS-K4a, Schwere II, rang 5; herausgelöste Fassung: keine)
% Zweitthema ausgeblendet (Spalten <feld>_einheit, pruefheft-p.md v0.9)
\begin{pfaufg}{nach~P10\\’23~\pkt{}~4a}{5.}{}
Die Gerade $f$ hat die Gleichung $f(x) = 3x - 2$. Zeichne $f$ in das Koordinatensystem. Entscheide, ob die Aussagen wahr oder falsch sind:
\par\smallskip\noindent\gpassend{\begin{ksys}[xmin=-2,xmax=3,ymin=-4,ymax=4,karo=0.35]\end{ksys}}\par
\gwahrfalsch{1}{Die Steigung von $f$ ist negativ.}\gwahrfalsch{2}{$f$ schneidet die y-Achse in $(-2\,|\,0)$.}
\fusshilfe{\mbox{5:~Gerade durch (0|$-$2) und (1|1); Aussagen 1 und 2 falsch}}
\end{pfaufg}

% ===================== Stufe: Gleichung ablesen oder aus Punkten aufstellen =====================
\gkopfvor{\pfstufekopf{Gleichung ablesen oder aus Punkten aufstellen}{selten \pkt{} \pfjahre{25}{19 \pkt{} 17 \pkt{} 15}}\gfrage{Die Steigung aus zwei Punkten: Wo steht die Formel? Formelsammlung, Seite \linie{12mm}}}
% 2025-OS-K5a-m1 (Original 2025-OS-K5a, Schwere II, rang 6; herausgelöste Fassung: keine)
% Zweitthema ausgeblendet (Spalten <feld>_einheit, pruefheft-p.md v0.9)
\begin{pfaufg}{nach~P10\\’25~\pkt{}~5a}{6.}{}
Der Graph einer linearen Funktion $f$ geht durch die Punkte $A(-2\,|\,5)$ und $B(3\,|\,-2{,}5)$. Zeichne den Graphen von $f$ in das Koordinatensystem. Entscheide, ob die Aussagen wahr oder falsch sind:
\par\smallskip\noindent\gpassend{\begin{ksys}[xmin=-3,xmax=4,ymin=-3,ymax=6,karo=0.35]\end{ksys}}\par
\gwahrfalsch{1}{Der Graph von $f$ steigt monoton.}\gwahrfalsch{2}{Der Graph von $f$ schneidet die y-Achse in $(0\,|\,2)$}
\par\smallskip\noindent Gib eine Gleichung von $f$ an.
\par\noindent\hfill y $=$\linie{40mm}
\fusshilfe{\mbox{6:~Gerade durch A und B; Aussage 1 falsch, Aussage 2 wahr; $-$1,5x + 2}}
\end{pfaufg}

% ===================== Stufe: Nullstelle =====================
\gkopfvor{\pfstufekopf{Nullstelle}{selten \pkt{} \pfjahre{22}{19}}}
% kein Einstieg: keine Schülerfrage in der Gliederung (Schülerfrage, -Fälle, -Form, -Weg, -Skizze; oder Frage) – Einstieg nur von dort, nie aus der Bank (pruefung.md § 5, 12.10.)
% 2022-OS-K3a-m1 (Original 2022-OS-K3a, Schwere I, rang 7; herausgelöste Fassung: keine)
% Zweitthema ausgeblendet (Spalten <feld>_einheit, pruefheft-p.md v0.9)
\begin{pfaufg}{nach~P10\\’22~\pkt{}~3a}{7.}{}
Gegeben ist die lineare Funktion $f(x) = -2x + 5$. Zeichne den Graphen von $f$ in das Koordinatensystem. Gib die Nullstelle von $f$ an.
\par\smallskip\noindent\gpassend{\begin{ksys}[xmin=-2,xmax=4,ymin=-2,ymax=6,karo=0.35]\end{ksys}}\par
\par\noindent\hfill x $=$\linie{30mm}
\fusshilfe{\mbox{7:~Gerade durch (0|5) und (1|3); 2,5}}
\end{pfaufg}

% ===================== Stufe: Punktprobe =====================
\gkopfvor{\pfstufekopf{Punktprobe}{immer \pkt{} \pfjahre{26 \pkt{} 25 \pkt{} 24 \pkt{} 23}{}}}
% kein Einstieg: keine Schülerfrage in der Gliederung (Schülerfrage, -Fälle, -Form, -Weg, -Skizze; oder Frage) – Einstieg nur von dort, nie aus der Bank (pruefung.md § 5, 12.10.)
% 2023-OS-B1i-m1 (Original 2023-OS-B1i, Schwere I, rang 8; herausgelöste Fassung: keine)
\begin{pfaufg}{’23~\pkt{}~B1i}{8.}{Basis~\pkt{}\\ohne TR}
Die Gerade $f$ hat die Gleichung $f(x) = -5x + 2$. Einer dieser Punkte liegt nicht auf $f$. Kreuze ihn an:
\par\smallskip\noindent\pfkreuz{A}{$(-2\,|\,12)$}\pfkreuz{B}{$(3\,|\,-12)$}\pfkreuz{C}{$(4\,|\,-18)$.}
\fusshilfe{\mbox{8:~B}}
\end{pfaufg}
% 2025-OS-B1h-m1 (Original 2025-OS-B1h, Schwere I, rang 9; herausgelöste Fassung: keine)
\begin{pfaufg}{’25~\pkt{}~B1h}{9.}{Basis~\pkt{}\\ohne TR}
Genau einer der Werte löst die Gleichung $x(x+4) = -3$. Kreuze ihn an:
\par\smallskip\noindent\pfkreuz{A}{$x = 1$}\pfkreuz{B}{$x = 3$}\pfkreuz{C}{$x = -2$}\pfkreuz{D}{$x = -3$.}
\fusshilfe{\mbox{9:~$-$3}}
\end{pfaufg}
% 2026-FOR-K5b-m1 (Original 2026-FOR-K5b, Schwere I, rang 10; herausgelöste Fassung: keine)
\begin{pfaufg}{nach~P10\\’26~\pkt{}~5b}{10.}{}
Prüfe mit einer Rechnung, ob der Punkt $P(-2\,|\,7)$ auf dem Graphen von $f(x) = -3x + 1$ liegt.
\karo{2}
\fusshilfe{\mbox{10:~7, P liegt auf dem Graphen}}
\end{pfaufg}
% 2024-OS-K3c-m1 (Original 2024-OS-K3c, Schwere I, rang 11; herausgelöste Fassung: keine)
% Zweitthema ausgeblendet (Spalten <feld>_einheit, pruefheft-p.md v0.9)
\begin{pfaufg}{nach~P10\\’24~\pkt{}~3c}{11.}{}
Die Parabel $p$ ist der Graph von $p(x) = x^2 - 5$. Prüfe mit einer Rechnung, ob der Punkt $T(-3\,|\,5)$ auf der Parabel $p$ liegt.
\karo{2}
\fusshilfe{\mbox{11:~Nein}}
\end{pfaufg}

% ===================== Stufe: Gleichung aus Text =====================
\Needspace*{0.16\textheight}\pfstufekopf{Gleichung aus Text}{oft \pkt{} \pfjahre{23 \pkt{} 22}{21 \pkt{} 16}}
% Schülerfrage der Stufe (Gliederung; Einstieg nur von dort, nie aus der Bank – 12.10.)
\begin{pfaufg}{}{12.}{}
Anfangswert oder Änderung je Schritt? Schreib zu jeder Zahl, was sie in der Gleichung y = m · x + n ist (m oder n). Rechne nicht.
\par\smallskip
\par\noindent\makebox[\mbpfteil][l]{\textbf{a)}}Handytarif: 8 € im Monat und 0,09 € je Minute\hfill\linie{30mm}
\par\noindent\makebox[\mbpfteil][l]{\textbf{b)}}Kerze: 25 cm hoch, brennt je Stunde 2 cm ab\hfill\linie{30mm}
\par\fusshilfe{\mbox{12:~a) 8 ist n, 0,09 ist m; b) 25 ist n, $-$2 ist m}}
\end{pfaufg}
% 2022-OS-K6a-m1 (Original 2022-OS-K6a, Schwere I, rang 12; herausgelöste Fassung: keine)
\begin{pfaufg}{nach~P10\\’22~\pkt{}~6a}{13.}{}
Maxi und Paula haben zum 16. Geburtstag Geld bekommen und legen jeden Monat etwas dazu, um den Führerschein zu bezahlen. Maxi beginnt mit 1 360 € und spart 24 € im Monat. Paula beginnt mit 820 € und spart 65 € im Monat. Nach einem Jahr hat Paula 1 600 € gespart. Berechne, wie viel Geld Maxi nach einem Jahr hat.
\karo{2}
\fusshilfe{\mbox{13:~1648 €}}
\end{pfaufg}
% 2022-OS-K6b-m1 (Original 2022-OS-K6b, Schwere II, rang 13; herausgelöste Fassung: keine)
% Vorspann wie Nr. 13 (2022-OS-K6a-m1), gekürzt
\begin{pfaufg}{nach~P10\\’22~\pkt{}~6b}{14.}{FOR}
Angaben wie in Nr. 13. Stelle eine Gleichung für Paulas Guthaben auf ($y$: Guthaben in €, $x$: Zahl der Monate). Für Fahrschule und Prüfung braucht Paula 1 800 €. Bestimme, wie viele Monate sie mindestens sparen muss.
\karo{3}
\fusshilfe{\mbox{14:~65x + 820; 16 Monate}}
\end{pfaufg}
% 2023-OS-K3a-m1 (Original 2023-OS-K3a, Schwere II, rang 14; herausgelöste Fassung: keine)
\begin{pfaufg}{nach~P10\\’23~\pkt{}~3a}{15.}{}
Herr Mert braucht für seinen Umzug einen Transporter. Angebot 1: 25 € Miete je Tag, 320 km insgesamt frei, jeder weitere Kilometer kostet 0,35 €. Angebot 2: einmalig 115 € Miete für eine Woche, jeder Kilometer kostet 0,08 €. Die Kosten für einen Miettag sollen für beide Angebote in einem Diagramm dargestellt werden. Gib zu jedem Angebot das passende Diagramm an.
\par\smallskip\noindent\setlength{\gbildbreite}{\dimexpr0.96\linewidth/4\relax}\sbox{\gbox}{\vbox{\begin{tikzpicture}[x=1.3cm,y=1.3cm,line width=0.6pt]\draw[-{Stealth[length=1.5mm]}] (0,0)--(1.15,0) node[below left,font=\tiny] {km};\draw[-{Stealth[length=1.5mm]}] (0,0)--(0,1.05) node[below right,font=\tiny] {€};\draw[line width=0.8pt] (0,0)--(1,0.85);\end{tikzpicture}\par\begin{tikzpicture}[x=1.3cm,y=1.3cm,line width=0.6pt]\draw[-{Stealth[length=1.5mm]}] (0,0)--(1.15,0) node[below left,font=\tiny] {km};\draw[-{Stealth[length=1.5mm]}] (0,0)--(0,1.05) node[below right,font=\tiny] {€};\draw[line width=0.8pt] (0,0.35)--(1,0.6);\end{tikzpicture}\par\begin{tikzpicture}[x=1.3cm,y=1.3cm,line width=0.6pt]\draw[-{Stealth[length=1.5mm]}] (0,0)--(1.15,0) node[below left,font=\tiny] {km};\draw[-{Stealth[length=1.5mm]}] (0,0)--(0,1.05) node[below right,font=\tiny] {€};\draw[line width=0.8pt] (0,0.3)--(0.45,0.3)--(1,0.85);\end{tikzpicture}\par\begin{tikzpicture}[x=1.3cm,y=1.3cm,line width=0.6pt]\draw[-{Stealth[length=1.5mm]}] (0,0)--(1.15,0) node[below left,font=\tiny] {km};\draw[-{Stealth[length=1.5mm]}] (0,0)--(0,1.05) node[below right,font=\tiny] {€};\draw[line width=0.8pt] (0,0.7)--(0.45,0.7)--(1,0.25);\end{tikzpicture}}}\gbildname{A}{\begin{tikzpicture}[x=1.3cm,y=1.3cm,line width=0.6pt]\draw[-{Stealth[length=1.5mm]}] (0,0)--(1.15,0) node[below left,font=\tiny] {km};\draw[-{Stealth[length=1.5mm]}] (0,0)--(0,1.05) node[below right,font=\tiny] {€};\draw[line width=0.8pt] (0,0)--(1,0.85);\end{tikzpicture}}{\dimexpr\ht\gbox/4+\baselineskip\relax}\hfill\gbildname{B}{\begin{tikzpicture}[x=1.3cm,y=1.3cm,line width=0.6pt]\draw[-{Stealth[length=1.5mm]}] (0,0)--(1.15,0) node[below left,font=\tiny] {km};\draw[-{Stealth[length=1.5mm]}] (0,0)--(0,1.05) node[below right,font=\tiny] {€};\draw[line width=0.8pt] (0,0.35)--(1,0.6);\end{tikzpicture}}{\dimexpr\ht\gbox/4+\baselineskip\relax}\hfill\gbildname{C}{\begin{tikzpicture}[x=1.3cm,y=1.3cm,line width=0.6pt]\draw[-{Stealth[length=1.5mm]}] (0,0)--(1.15,0) node[below left,font=\tiny] {km};\draw[-{Stealth[length=1.5mm]}] (0,0)--(0,1.05) node[below right,font=\tiny] {€};\draw[line width=0.8pt] (0,0.3)--(0.45,0.3)--(1,0.85);\end{tikzpicture}}{\dimexpr\ht\gbox/4+\baselineskip\relax}\hfill\gbildname{D}{\begin{tikzpicture}[x=1.3cm,y=1.3cm,line width=0.6pt]\draw[-{Stealth[length=1.5mm]}] (0,0)--(1.15,0) node[below left,font=\tiny] {km};\draw[-{Stealth[length=1.5mm]}] (0,0)--(0,1.05) node[below right,font=\tiny] {€};\draw[line width=0.8pt] (0,0.7)--(0.45,0.7)--(1,0.25);\end{tikzpicture}}{\dimexpr\ht\gbox/4+\baselineskip\relax}\par
\par\noindent\hfill Angebot 1 $=$\linie{30mm}
\par\noindent\hfill Angebot 2 $=$\linie{30mm}
\fusshilfe{\mbox{15:~C; B}}
\end{pfaufg}
% 2023-OS-K3b-m1 (Original 2023-OS-K3b, Schwere II, rang 15; herausgelöste Fassung: keine)
% Vorspann wie Nr. 15 (2023-OS-K3a-m1), gekürzt
\begin{pfaufg}{nach~P10\\’23~\pkt{}~3b}{16.}{FOR}
Angaben wie in Nr. 15. Herr Mert mietet den Transporter von Montag bis Freitag und fährt insgesamt 450 km. Berechne die Kosten für beide Angebote. Gib an, welches billiger ist. Stelle für Angebot 2 eine Gleichung auf, mit der man die Kosten einer Woche für jede Kilometerzahl berechnen kann.
\karo{4}
\fusshilfe{\mbox{16:~Angebot 1: 170,50 €; Angebot 2: 151,00 €; Angebot 2 günstiger; 0,08x + 115}}
\end{pfaufg}

\pfblattende{11}{26 FOR & 25 & 24 & 23 & 22 & 21 & 20 & 19 & 18 & 17 & 16}{5a 5b & 1f 1h 5a & 1i 3a 3c & 1i 3a 3b 4a & 3a 6a 6b & 2a 2b 6a 6c 6d 7a &  & 2a 2b 2c &  & 5a 5b & 1c 6a 6b 6c}%
{\pfausschnitt{}{Lineare Funktionen}{}}
\end{document}